% Professional Two-Column Research Paper Format
\documentclass[conference,twocolumn]{IEEEtran}

% Essential packages
\usepackage{amsmath,amssymb}
\usepackage[T1]{fontenc}
\usepackage[utf8]{inputenc}
\usepackage{textcomp}
\usepackage{xcolor}
\usepackage{booktabs,array}
\usepackage{multirow}
\usepackage{tabularx}
\usepackage{graphicx}
\usepackage{cite}
\usepackage{hyperref}
\usepackage{url}
\usepackage{balance}
\usepackage{enumitem}

% For tables that span both columns
\usepackage{stfloats}

% Configure hyperref
\hypersetup{
  pdftitle={Design and Development of an LLM-Powered Document Parsing Pipeline for Automated Medical Data Extraction and Validation},
  pdfauthor={Dhruva Singh, Dr. Kanika Garg},
  hidelinks,
  colorlinks=false
}

% Tight list formatting
\providecommand{\tightlist}{%
  \setlength{\itemsep}{0pt}\setlength{\parskip}{0pt}}

% Set default figure placement
\makeatletter
\def\fps@figure{htbp}
\makeatother

% Prevent overfull lines
\setlength{\emergencystretch}{3em}

% Make section headings bold
\renewcommand{\thesection}{\Roman{section}}
\renewcommand{\thesubsection}{\Alph{subsection}}

\makeatletter
\renewcommand{\section}{\@startsection{section}{1}{\z@}%
  {-3.5ex \@plus -1ex \@minus -.2ex}%
  {2.3ex \@plus.2ex}%
  {\normalfont\large\bfseries\centering}}
\renewcommand{\subsection}{\@startsection{subsection}{2}{\z@}%
  {-3.25ex\@plus -1ex \@minus -.2ex}%
  {1.5ex \@plus .2ex}%
  {\normalfont\normalsize\bfseries\itshape}}
\makeatother

\begin{document}

% Paper Title
\title{Design and Development of an LLM-Powered Document Parsing Pipeline for Automated Medical Data Extraction and Validation}

% Author Information
\author{
\IEEEauthorblockN{Dhruva Singh}
\IEEEauthorblockA{Dept. of Computer Science and Engineering\\
SRMIST Delhi-NCR Campus\\
Ghaziabad, India\\
singhdhruva45@gmail.com\\
\url{https://orcid.org/0009-0000-3213-9187}}
\and
\IEEEauthorblockN{Dr. Kanika Garg}
\IEEEauthorblockA{Dept. of Computer Science and Engineering\\
SRMIST Delhi-NCR Campus\\
Ghaziabad, India\\
kanikap@srmist.edu.in}
}

\maketitle

% Abstract
\begin{abstract}
\textbf{Background:} Processing medical documents by hand, especially semi-structured documents like medical bills for healthcare reimbursements, presents a major bottleneck characterized by exorbitant costs, long waits, and high rates of human error. The complexity is compounded by the diverse scripts and variable writing styles common in Indian languages and the challenges of deciphering non-standardized medical records. This work details the design, development, and evaluation of a document parsing pipeline utilizing a Large Language Model (LLM) for automatically extracting and validating data from Indian medical bills. We implement a nuanced multi-stage approach to solve the problem of unstructured, variable-layout documents.

\textbf{Method:} The pipeline is constructed on a modified ensemble detection model consisting of a Weighted Box Fusion (WBF) of two deep learning architectures, YOLOv11n and Faster RCNN, to accurately and precisely detect regions. This use of YOLO-based detection is supported by its proven high accuracy in Region of Interest (ROI) localization on physical document images. The Optical Character Recognition (OCR) step extracts text after detection of the targeted regions' coordinates, which uses PaddleOCR as the primary engine and Tesseract as a fallback. Complicated tabular data, which requires Line-Item Recognition (LIR), is parsed using the Microsoft Table-Transformer model. The next step is the post-processing loop, which employs Google's MedGemma 4B, a 4-bit quantized multimodal model, for domain-specific validation, semantic correction, and intelligent text parsing, addressing the unique complexities of medical data. The entire pipeline runs over an async FastAPI, and the extracted text records are stored in a serverless PostgreSQL database.

\textbf{Results:} The ensemble detection model achieved a Macro F1-score of 0.906 at a 0.5 IoU threshold. The pipeline is capable of processing a single document in approximately 15--20 seconds; this latency reflects the priority placed on high-fidelity domain-specific validation via the multi-stage approach and specialized LLM, contrasting with faster, purely OCR-free models.

\textbf{Conclusion:} Utilizing the pipeline, a medical bill can be parsed, extracting the required fields and storing them after validation in a serverless database. This approach successfully combines advanced visual recognition and domain-aware language modelling to tackle information extraction from complex, unstructured documents.
\end{abstract}

\begin{IEEEkeywords}
Medical bills, OCR, Ensemble model, Object detection, YOLOv11, Faster R-CNN, MedGemma, Extraction, Validation, Document parsing
\end{IEEEkeywords}

\section{Introduction}

The healthcare industry faces an enormous challenge of processing
medical bills and insurance claims, which is traditionally undergone via
manual and labour-intensive techniques \cite{ref1}. This results in possible
errors and high costs, sometimes leading to denials and rejections of
claims and reimbursements after an overall lengthy procedure \cite{ref1}.
Automated document parsing pipelines can minimize this manual burden and
perform more efficiently by extracting structured relevant records from
semi or unstructured formats \cite{ref1}, \cite{ref2}. However, healthcare
documents are highly variable in their formats and noisy images/scans
can limit normal OCR based pipelines \cite{ref3}, \cite{ref4}. I tried to
address this by designing and implementing an end-to-end OCR-LLM based
pipeline \cite{ref3}, \cite{ref4}, \cite{ref5} that can understand context, detect
positioning of targeted fields, extract those relevant fields and
finally store them in a structured, uniform format in a serverless
database. Some key features of the proposed pipeline include (1) an
ensemble model to detect the regions with the fields of interest; (2) a
dual-engine OCR and table-parsing for text extraction; (3) an LLM based
layer for sematic validation and correction of outputs; and finally (4)
the ability to run locally on lightweight hardware. This paper
highlights each step, ranging from data collection, preprocessing,
models' training, OCR and LLM integration to storing results in a
serverless database and how the stated implementation can significantly
improve the efficiency over the current legacy methods.

\section{Related Works}

\subsection{Traditional and OCR Based Document Understanding}

Most of the document parsing attempts in the past have been based on
handcrafted features and rule-based methods for text detection and
extraction, which generally have shown very low-quality performance on
noisy or unstructured layouts \cite{ref3}, \cite{ref5}. In earlier days,
document understanding approaches had been either probabilistic,
rule-based, or feature-based with the features being manually designed,
which almost always failed in the case of unfamiliar templates \cite{ref5}.
Standard Optical Character Recognition (OCR) methods made their way from
rule-based and statistical methods like Hidden Markov Models (HMMs) and
Support Vector Machines (SVMs) to deep learning techniques, chiefly
employing LSTMs (a kind of Recurrent Neural Network) and
Transformer-based models such as TrOCR \cite{ref5}. Programs like Tesseract,
EasyOCR, and MMOCR employ CNNs, RNNs, and the attention mechanism to
achieve foolproof script recognition \cite{ref2}, \cite{ref5}.

Although deep learning-based models have enhanced performance on printed
texts, they usually do not have sufficient adaptability to handle
variably positioned fields or multilingual contents, particularly in
unstructured documents like medical bills \cite{ref4}, \cite{ref5}. For example,
the creation of a strong OCR system for Indian languages is naturally
tough because of the variety of scripts, intricacy of characters, and
multilingual text format, which in turn complicates traditional OCR
methods \cite{ref5}. In addition, handwriting recognition is still a very
challenging problem mainly because of the variability in writing style
and the fact that documents like medical prescriptions are usually
non-standard in their layout \cite{ref6}.

Several neural network-based parsers such as CRAFT \cite{ref7} and PixelLink
\cite{ref8}, have enhanced the spatial localization of text in the documents
that have been scanned and have made it possible to represent the texts
as polygons and curves. Also, the tools which use both visual and
positional information like LayoutParser \cite{ref9} and the LayoutLM series
\cite{ref10} have enhanced the semantic understanding of the structure of
the visually-rich documents by combining the textual, visual, and
spatial cues.

Nevertheless, these models are in general designed as two-stage
pipelines \cite{ref3}, \cite{ref11} which indicates that they are still greatly
dependent on the quality of the input provided by an external OCR system
\cite{ref4}. Such reliance on external OCR restricts the performance of the
system in the case of low-quality scans or handwritten documents
\cite{ref3}, \cite{ref4}, \cite{ref11}. The dependence on external OCR might result
in the drop of performance if there is a change in the domain, and the
errors made by the OCR engine can be carried forward and have a negative
effect on the following document understanding processes \cite{ref4}.
\subsection{Hybrid Visual-Semantic Approaches}

To identify key regions in documents, several studies delved into deep
object detection models as a way to go beyond the limitations of
template-based methods \cite{ref9}, \cite{ref12}. The object detection methods
such as Faster R-CNN \cite{ref13} and Mask R-CNN \cite{ref14} are largely
applied in the setting of documents for the purpose of locating the
structural elements like tables and figures. While the Faster
R-CNN-based solutions were able to deliver higher accuracy with a slower
processing time, their YOLO-based equivalents (such as YOLOv7, YOLOv8,
and YOLOv5) were characterized by a very fast inference speed \cite{ref6},
\cite{ref15}. In the case of current medical scenarios, YOLO-based models
are utilized mainly for Region of Interest (ROI) detection in
prescriptions, thereby achieving a very high accuracy (99.6\%) in
localizing medicine-related data \cite{ref6}.

In particular, for variable-layout forms, the method of CRAFT (Character
Region Awareness for Text Detection) brought improvement to spatial
localization by examining not only individual character regions but also
the affinities between characters, thus being highly adaptable to
arbitrarily-oriented or curved texts \cite{ref5}, \cite{ref7}. More complex
pipelines have introduced dual-stage OCR processing with
confidence-based fallback strategies and domain-specific layout analysis
through table parsers \cite{ref9}. The degree to which the use of dedicated
table parsing models is effective can be gauged by the performance of
benchmarks like DocILE that treat DETR-based Table Transformer as a
baseline for complex Line-Item Recognition (LIR) tasks \cite{ref2}. Besides
that, the Docling toolkit has state-of-the-art highly specialized AI
models such as the layout analysis model (based on RT-DETR trained on
DocLayNet) and TableFormer for table structure recognition that can
convert complex documents into unified, structured formats \cite{ref16}.
TableFormer is mainly targeted at dealing with different types of
problems in tables such as cell spans, hierarchy, and inconsistency of
alignment \cite{ref16}.

Though these systems excelled on structured documents, they frequently
failed to generalize well across various formats, such as Indian medical
bills, because of limitations arising from script variability,
multilingual content, and noisy backgrounds. The recognition of
documents from developing countries such as medical prescriptions is
inherently a tough task since they are largely handwritten and physical
and usually have dirty writing styles and non-standardized formats
\cite{ref6}. The complexity is magnified by the fact that in places like
India there are different scripts, complicated characters, ligatures,
and diacritics. Moreover, the use of mixed languages (code-mixing) is
the continual problem of traditional OCR systems in these areas. The
need of handling complex scripts, multilingual content, and noisy inputs
is the reason behind the trend of switching to advanced Vision-Language
Models (VLMs) which combine visual and linguistic cues for a deeper
understanding of text and are quite often on par with or even better
than traditional OCR in challenging environments \cite{ref5}.

\subsection{LLM-Driven Post-OCR Validation}

Firstly, VLMs and LLMs have been utilized to correct OCR and validate the
semantics of the text \cite{ref5}. Latest improvements in Vision-Language
Models (VLMs) are to use both visual and language clues for a deeper
understanding of the text. They often outperform traditional OCR systems
because they use language priors to understand even ambiguous or
deteriorated texts \cite{ref5}. Due to this, these systems can process an
entire page in a more holistic way and the need for explicit layout
detection and word cropping is getting less and less \cite{ref5}.

The examples are Donut, Pix2Struct, and Dessurt that demonstrate that by
adding layout reasoning and contextual understanding, the outputs can be
more structured thus there is often no need for traditional OCR to be
done \cite{ref4}, \cite{ref5}, \cite{ref11}, \cite{ref17}. Donut, an OCR-free document
understanding transformer, uses an end-to-end architecture (a Swin
encoder and BART-like decoder) to directly connect the input image with
the desired structured output sequence \cite{ref4}, \cite{ref5}. Likewise,
Dessurt is a new end-to-end document understanding transformer that,
along with text recognition, also performs document understanding,
thereby producing arbitrary text in an autoregressive manner as output
when given an image and a task string as inputs \cite{ref11}.

Still, these networks are very demanding in terms of computing power and
require end-to-end fine-tuning, which

may not be always possible for domain-specific, low-resource datasets
\cite{ref3}. For example, Donut is very heavy on the pretraining part and
thus requires a lot of pretraining steps on large datasets such as
IIT-CDIP and may need up to 64 A100 GPUs at the same time for
pretraining steps \cite{ref4}, \cite{ref18}. Though scaling VLMs results in
better performance, it leads to larger models with increased
computational complexity. Some OCR-free methods like DocParser even
suggest that end-to-end approaches can be up to two times slower than
optimized state-of-the-art OCR-dependent methods, thus demonstrating the
price that is paid in terms of computation when making the trade-off
\cite{ref3}. The next generation of VLMs is actively exploring ways to
reduce their computational load by means of quantization and pruning so
that they can be more easily deployed on mobile and embedded devices
\cite{ref5}.

The Indian medical context has been the subject of previous works and a
generalized Indic OCR has likewise been the focus of attention for
researchers working on the above-mentioned topics \cite{ref5}. Special
emphasis has been put on the Indian languages' models
and the use of script-sensitive tokenizers for supporting Hindi and
regional OCR tasks. The research team has adapted the HindiOCR-VLM model
for Hindi OCR by utilizing parameter-efficient fine-tuning (LoRA) on
models pre-trained for Chinese and English, thus showing the
model's capability with the Devanagari script \cite{ref5}.

Most of such domain-specific models are either layout
predictor-dependent or continue being domain-specific that limits their
usage across different document types, however \cite{ref5}. Besides that,
traditional IE pipelines in the healthcare sector, particularly in
developing countries, are failing due to the very nature of
non-standardized formats, messy writing styles, and specialized
terminologies typically found in handwritten and physical prescriptions
\cite{ref6}. The most common specialized systems may attain high text
extraction performance but, at the same time, they might not preserve
the context or relationship among terms (e.g., dosage, instructions, and
medicine name) \cite{ref6}.

These trends have been complemented by my strategy of presenting a
lightweight, rule-guided pre-validation layer followed by a quantized,
multimodal LLM fine-tuned for the medical domain, MedGemma \cite{ref3},
\cite{ref19}. The intent of this layout is to lessen the dependence on the
fully end-to-end training and to detour the heavy load that large VLMs
bring with them. In contrast to general-purpose models, MedGemma makes
it possible for semantic correction, contextual reformatting, and the
validation of OCR-extracted data to be carried out, thus the reliability
is upgraded without the necessity of an extensive retraining \cite{ref2}.
The mentioned plan leads to a solution that integrates robust OCR
processing with advanced language models for taking on complex tasks
such as document information extraction (IE), which entails the
comprehension of semantics, layout, and context. Moreover, by employing
serving efficiency methods akin to virtual memory and paging (like
vLLM's PagedAttention) it would be feasible to guarantee
that a quantized LLM like MedGemma is operating in a lightweight manner,
thereby allowing maximum throughput even in situations where resources are limited \cite{ref19}, \cite{ref20}.

\section{Data Collection and Detection Models}

\subsection{Data Collection, Labeling, Augmentation and Splitting}

I came across a publicly accessible dataset containing 395 Indian medical
bills with some basic annotations on the Roboflow Universe, but those
annotations were not suitable for my purpose. The area of medical
prescriptions and bills, especially for developing countries, is very
challenging because these documents are mainly in hard copies and are
not in a standardized format \cite{ref6}. In addition, the text can be in
different languages and use different scripts, which makes the job of an
OCR even more difficult \cite{ref5}.

I forked the dataset to my workspace and changed the annotation schema
to have eight classes: (1) \textit{date\_of\_receipt}, (2) \textit{gstin},
(3) \textit{invoice\_no}, (4) \textit{mobile\_no}, (5)
\textit{product\_table}, (6) \textit{store\_address}, (7) \textit{store\_name}
and (8) \textit{total\_amount}. The emphasis on locating and extracting
the main fields, in particular, the \textit{product\_table} (line items),
is in perfect harmony with the described complex tasks of Key
Information Localization and Extraction (KILE) and Line-Item Recognition
(LIR) in document processing benchmarks \cite{ref2}. I then re-annotated the
images of the newly defined labels through Roboflow' s
bounding-box annotation tool.

In order to increase generalization, I employed two data-augmentation
techniques: (1) image rotation (clockwise and anti-clockwise by 90°) to
stimulate orientation variance \cite{ref8}. This augmentation technique is
generally applied in the training of object detection models like YOLO
and in Vision-Language Models (VLMs) \cite{ref15}. And (2) grayscale
filtering was performed on the randomly chosen 15\% of the images to
represent a badly scanned or faded printed document. The idea of adding
synthetic noise and blur or changing brightness and contrast is a
perfect solution to improve the models' robustness and
generalization abilities to poor-quality or degraded document images
\cite{ref3}, \cite{ref5}, \cite{ref6}.

\begin{table}[htbp]
\caption{Dataset Class Distribution Across Train, Validation, and Test Sets}
\label{tab:dataset}
\centering
\footnotesize
\begin{tabular}{lcccc}
\toprule
Class & Train & Valid & Test & Total \\
\midrule
date\_of\_reciept & 724 & 59 & 58 & 841 \\
gstin & 716 & 58 & 57 & 831 \\
invoice\_no & 720 & 59 & 58 & 837 \\
mobile\_no & 679 & 60 & 57 & 796 \\
product\_table & 815 & 60 & 68 & 943 \\
store\_address & 710 & 59 & 58 & 827 \\
store\_name & 839 & 65 & 68 & 972 \\
total\_amount & 701 & 57 & 58 & 816 \\
\midrule
Total & \textbf{5904} & \textbf{477} & \textbf{482} & \textbf{6863} \\
\bottomrule
\end{tabular}
\end{table}

Consequently, after performing these augmentations, I had a larger
dataset of 862 images, which was divided into training: 742 images
($\sim$86.1\%), validation: 60 images ($\sim$7\%), and testing: 60 images ($\sim$7\%)
sets. Eventually, this dataset was converted into YOLOv11 and COCO
formats for the purpose of training two different object detection
architectures.

\subsection{YOLOv11n Model Training}

Initially, I trained a YOLOv11n (nano variant) \cite{ref15} object detection
model. Training ran for 100 epochs with a batch size of 16 on
CUDA-enabled GPU with an image input resolution of 640×640 pixels, and
the patience of early stopping was set to 50 epochs for preventing
overfitting. The architecture of YOLOv11n consists of an efficient
nano-scale backbone optimized for mobile and edge deployment with
scaling factors of depth and width. CSPDarknet-inspired feature
extraction in the form of multi-scale detection with PANet-style feature
pyramid networks uses anchor-free decoding with distribution focal loss
(DFL) for accurate bounding box regression.

The optimization strategy utilized stochastic gradient descent with a
cosine annealing learning rate schedule, initialized at 0.0008276,
linearly decaying to 0.0000166 by epoch 100, applied uniformly across
all three parameter groups: pg0, pg1, pg2. The training pipeline
utilized comprehensive data augmentation strategies: (1) mosaic
augmentation of four images, (2) mixup through blending image pairs, (3)
HSV color space transformations in hue, saturation and value jittering,
(4) random horizontal flips, (5) scale jittering, (6) translation
augmentation, and (7) affine transformation to improve model
generalization and robustness to lighting variation, orientation change,
and document quality degradation. The multi-component loss function
combined box regression loss based on CIoU for precise localization,
classification loss as binary cross-entropy with focal loss weighting to
balance class imbalance, and DFL for refining the corners of bounding
boxes in a finer granularity, with final training losses converging to
1.028 for box, 0.703 for class, and 0.924 for DFL.

\begin{table}[htbp]
\caption{YOLOv11n Model Performance Metrics}
\label{tab:yolo}
\centering
\footnotesize
\begin{tabular}{lccc}
\toprule
Metric & Train & Valid & Test \\
\midrule
Precision & 0.8876 & 0.8925 & 0.9023 \\
Recall & 0.8793 & 0.8826 & 0.8825 \\
F1 & 0.8834 & 0.8865 & 0.8923 \\
mAP@0.5 & 0.9073 & 0.9076 & 0.9217 \\
mAP@0.5-0.95 & 0.6070 & 0.6107 & 0.6091 \\
\bottomrule
\end{tabular}
\end{table}

Training convergence saw steady improvement across all metrics, with
precision increasing from 0.19\% in epoch 1 to 88.76\% in epoch 100,
recall improving from 11.64\% to 87.93\%, and mAP50 rising from 4.76\%
to 90.73\%, whereas mAP50-95 was 60.70\%, showing excellent performance
at the multi-threshold. Convergence in validation losses went down from
starting values of 2.353 for box, 3.947 for class, and 1.420 for DFL to
finishing converged values of 1.092, 0.674, and 0.933, respectively. On
this test set alone, the model attained 90.23\% precision, 88.25\%
recall, and 92.17\% mAP50 with exceptionally high performance on
\textit{product\_tables} and \textit{total\_amount} at 99.5\% mAP50/100\%
recall and 99.48\% mAP50/100\% recall, respectively, while
\textit{mobile\_no} detection presented difficulties, 64.24\% mAP50 and
49.12\% recall, because of their varied formatting and the complexity of
OCR-style recognition. The weights had been saved on epoch 88, where the
validation mAP50 value peaked at 90.87\%, proving how very well YOLOv11n
was suited for real-time document understanding tasks with very low
computational overhead and, correspondingly, ready-for-deployment
features.

\subsection{Faster R-CNN (ResNet-50 FPN) Model Training}

Next, I trained a Faster R-CNN \cite{ref13}, \cite{ref14} object detection
model, using the ResNet-50 Feature Pyramid Network backbone, in PyTorch
TorchVision. The model underwent training for 25 epochs with a batch
size of 4 on CUDA-enabled GPU using COCO-pretrained ResNet-50 FPN
weights by way of transfer learning, whereby pre-learned hierarchical
feature representations from natural images greatly accelerated model
convergence and improved generalization to document-domain object
detection.

The Faster R-CNN architecture adopts a two-stage detection pipeline
wherein a class-agnostic region proposal is generated through
anchor-based sliding window mechanisms by the RPN, followed by a
subsequent classification and refinement of the bounding box through the
ROI head. A ResNet-50 FPN backbone extracts multi-scale feature pyramids
using a bottom-up pathway-a ResNet convolution with residual
connections-followed by a top-down pathway through lateral connections
with upsampling, allowing it to detect objects across varied scale
ranges, which is critical in the case of medical bills, where field
sizes range from small mobile numbers to large product tables. The model
architecture embeds four kinds of specialized losses: (1) RPN
classification loss-a binary cross-entropy-to provide objectness scores;
(2) RPN bounding box regression loss-a smooth L1-for the RPN to arrive
at proper proposal localizations; (3) ROI classifier loss-a
cross-entropy-to calculate final class predictions; and (4) ROI box
regression loss-smooth L1-for obtaining very accurate coordinate
refinements.

\begin{table}[htbp]
\caption{Faster R-CNN Model Performance Metrics}
\label{tab:fasterrcnn}
\centering
\footnotesize
\begin{tabular}{lccc}
\toprule
Metric & Train & Valid & Test \\
\midrule
Precision & 0.8339 & 0.7857 & 0.7861 \\
Recall & 0.8747 & 0.7919 & 0.7989 \\
F1 & 0.8529 & 0.7882 & 0.7917 \\
mAP@0.5 & 0.9266 & 0.9043 & 0.8827 \\
\bottomrule
\end{tabular}
\end{table}

The optimization strategy used SGD with momentum of 0.9 and L2 weight
decay regularization of 0.0005 to prevent overfitting, starting from an
initial learning rate of 0.005. Furthermore, the step-based learning
rate schedule was applied with the step size of 10 epochs and decay
factor (gamma) of 0.1, reducing the learning rate to 0.0005 at epoch 10
and further to 0.00005 at epoch 20 to enable fine-grained convergence.
Training convergence showed a rapid initial improvement followed by
steady refinement: the total loss decreased from 0.931 (epoch 1) to
0.149 (epoch 25), and the component losses stabilized at 0.0285
(classifier), 0.0669 (box regression), 0.0074 (objectness), and 0.0462
(RPN box regression). The loss decomposition shows that RPN components
converged faster than the ROI head losses, which means effective
proposal generation while classification and localization required
extended fine-tuning.

The final evaluation metrics showed strong performance with mAP50 of
88.27\%, mAP75 of 66.92\%, and mAP50-95 of 58.67\%, which indicated the
robustness in the accuracy of object localization. The best weights of
the model were saved at epoch number 24, showing its suitability for
high-accuracy document understanding tasks where precise multi-scale
object localizations are needed.

\subsection{Weighted Box Fusion Ensemble Model}

After training the two models I created an ensemble model that combines
the Faster R-CNN with YOLOv11n, using Weighted Box Fusion to take
advantage of the complementary strengths of both architectures in
improving the accuracy of medical bill field detection. The ensemble
strategy followed addressed the fundamental trade-off between the two
base models: Faster R-CNN excels in precise localization, yielding
higher accuracy for structured fields such as \textit{gstin} and
\textit{store\_name} (88.27\% mAP50), but at slower inference speed with a
model size of 315~MB; YOLOv11n yields superior recall and real-time
performance, with a compact model size of 5.2~MB (92.17\% mAP50), which
sometimes gives less precise bounding boxes. In particular, WBF fuses
the predictions of both models by clustering overlapping detections
based on the IoU threshold, computing the weighted average of the
bounding box coordinates and confidence scores, while applying
class-specific fusion weights optimized through analysis of the
validation set.

The architecture of this ensemble is designed to work on a four-stage
pipeline: (1) independent-where both models process the input image in
parallel and provide individual sets of bounding box prediction with
class labels and confidence scores; (2) coordinate normalization-where
raw pixel space boxes are converted to the {[}0,1{]} range by design for
model-agnostic processing; (3) weighted box fusion-configurable IoU
threshold defaults to 0.55, and for skip-box threshold, 0.0001 WEIGHTS
default distributions per class favor YOLOv11n for \textit{product\_table}
at 58\% and \textit{total\_amount} at 70\%, while Faster R-CNN has more
weigh for \textit{gstin} at 55\% and \textit{store\_name} at 52\%; and (4)
confidence filtering-based on class thresholds to remove low-confidence
detections and reduce false positives.

\begin{table}[htbp]
\caption{WBF Ensemble Model Performance at Different IoU Thresholds}
\label{tab:wbf}
\centering
\footnotesize
\begin{tabular}{lccc}
\toprule
Metric & IoU@0.50 & IoU@0.75 & IoU@0.95 \\
\midrule
Precision & 0.8856 & 0.7673 & 0.0473 \\
Recall & 0.9315 & 0.8071 & 0.0498 \\
F1-Score & 0.9080 & 0.7867 & 0.0485 \\
Macro Precision & 0.8869 & 0.7672 & 0.0448 \\
Macro Recall & 0.9289 & 0.8018 & 0.0467 \\
Macro F1 & 0.9066 & 0.7835 & 0.0457 \\
Mean IoU & 0.8377 & 0.8594 & 0.9613 \\
Detection Acc. & 0.8315 & 0.6483 & 0.0249 \\
\bottomrule
\end{tabular}
\end{table}

Implementation follows the ensemble-boxes library in Python with
integrated PyTorch and Ultralytics YOLO frameworks; inference is managed
through wrapper classes-FasterRCNNWrapper and YOLOWrapper-that handle
model loading, prediction formatting, and coordinate transformations.
Performance on the 60-image test set was significantly improved: at IoU
0.5, the ensemble reached an F1 score of 90.80\% compared to 89.23\% and
79.17\% baselines given by YOLO and Faster R-CNN, respectively, with the
best recall of 93.15\% and precision of 88.63\%. Per-class analysis
showed that fusion effectively reduced false positives when detecting
\textit{mobile\_no} and improved boundary precision for
\textit{product\_table} of variable size. The final ensemble model
requires both base model weights, 320 MB in total and provides a
well-balanced trade-off between accuracy and flexibility for deployment
in the document processing pipeline.

\section{Pipeline Architecture}

The architecture is designed as a modular, 6-step, high-performance
pipeline orchestrated by a central API service. At the high level, the
architecture starts with the user who interacts with the system via a
POST request to one of the FastAPI endpoints, which gets documented and
exposed via Swagger UI. This request contains the medical bill image and
triggers the 6-step processing workflow.

One of the main constituents of the architecture is the Eager Model
Loading mechanism. At the time the FastAPI application fires up. It
pre-loads into memory all five core models: YOLOv11n, Faster R-CNN,
PaddleOCR, Table-Transformer, and MedGemma. This is to ensure that all
computational resources are ready to go for inference and will not
introduce any cold-start latency for the first user request or
subsequent ones. The FastAPI service interfaces with the Neon PostgreSQL
database for job record creation and updates, the persistence of final
structured data, and state management for the processes.

\subsection{Stage-1: Document Upload and Initialization}

The pipeline starts with the upload of a medical bill document via the
POST \textit{/api/medical-bills/process} endpoint. The system supports
three image formats: JPEG, PNG, and PDF. Format detection is done
automatically, and format validation is performed. Upon receiving the
upload, the system performs initial quality checks: file size limits (up
to 10 MB), verification of the format by inspecting magic numbers, and
basic validation of image integrity. In the PostgreSQL database,
instantiate a unique \textit{MedicalBill} record with an auto-incremented
primary key identifier; save metadata like original \textit{filename},
\textit{file\_path}, \textit{file\_size} (in bytes), and
\textit{upload\_timestamp} (in microseconds). Set the document
\textit{status} to UPLOADED to lay the foundation for complete end-to-end
tracking in this pipeline. Store the uploaded file in the
system's temporary directory using
Python's tempfile module. Generate a random filename to
avoid conflicts. Create a new record for \textit{ProcessingJob} with
\textit{job\_type} as ``full\_pipeline''
and \textit{status} as ``PROCESSING'', which
will enable detailed tracking of the processing workflow.

Thus, a comprehensive set of error handling is implemented in this
initialization phase by making use of try-catch blocks to catch the
upload failure, database connection error, or some file system issue,
through which the document \textit{status} will be automated to FAILED and
details of errors will be logged for further debugging purposes. The
average time this stage takes is around 0.2-0.5 seconds, which is less
than 3\% of the average time taken by a document to go through the
entire processing pipeline.

\subsection{Stage 2: Ensemble-Based Region Detection}

After successful initialization, the document moves to the region
detection stage, which uses the WBF ensemble model combining YOLOv11n
and Faster R-CNN. The system reads the uploaded image with
OpenCV's function \textit{cv2.imread()}, returning a NumPy
array representation with shape (H, W, 3) in BGR color space. A new
\textit{ProcessingJob} record is created to track this stage independently
with \textit{job\_type} as
``ensemble\_detection''.

Both detection models process the image parallelly on a CUDA-enabled
GPU: YOLOv11n performs the inference at 640×640 resolution using its
anchor-free detection head with distribution focal loss, while the
Faster R-CNN processes the image through its ResNet-50 FPN backbone with
region proposal network and ROI head. Prediction set returned by each
model is a list of tuples of the form (\textit{bbox}, \textit{class\_label},
\textit{confidence\_score}), with bbox coordinates in {[}x1, y1, x2, y2{]}
format representing top-left and bottom-right corners. WBF normalizes
the coordinates in a {[}0, 1{]} range, clusters the overlapping
detections with IoU Threshold set to 0.55, and calculates weighted
averages of bounding box coordinates as well as confidence scores based
on class-specific fusion weights.

Ensemble typically detects 7-9 regions in a document corresponding to
the eight classes: \textit{date\_of\_receipt}, \textit{gstin},
\textit{invoice\_no}, \textit{mobile\_no}, \textit{product\_table},
\textit{store\_address}, \textit{store\_name}, and \textit{total\_amount}.
Often, the same field might get detected multiple times when both models
predict it with great confidence, which gets resolved by WBF due to its
weighted averaging. The final ensemble predictions filtered with a
class-specific confidence threshold of 0.5 (default) to remove
low-confidence false positives. Detection metadata, including the
\textit{ensemble\_model\_used} as ``Faster-RCNN + YOLO11n WBF'', and the average
\textit{detection\_confidence} across regions gets persisted to the
database. This stage finishes in 2-3 seconds on a CUDA-enabled RTX 3050
GPU, which accounts for about 15\% of the total processing time. On
successful detection, the \textit{status} of the \textit{ProcessingJob} is
changed to ``COMPLETED'', and on returning
less than three detected regions by the ensemble (indication of an
invalid or severely degraded input image), it is marked as
``FAILED''.

\begin{figure}[htbp]
\centering
\includegraphics[width=\columnwidth]{image2.png}
\caption{Ensemble Detection Process}
\label{fig:ensemble}
\end{figure}

\begin{figure*}[t]
\centering
\includegraphics[width=\textwidth]{image1.png}
\caption{Pipeline Architecture Overview}
\label{fig:pipeline}
\end{figure*}

\subsection{Dual-Engine OCR Extraction}

The third phase extracts the text in each of the recognized regions by a
complex dual-engine OCR method. To keep track of this phase, a
\textit{ProcessingJob} with \textit{job\_type} as
``ocr\_extraction'' is created. For each
bounding box given by the ensemble model, the program gets the
corresponding portion of the full image with NumPy array slicing plus a
5-pixel padding buffer. This extension is aimed at including the
characters, which may be a little outside the detected boundaries. The
cropped image, which is generally between 100×50 and 800×200 pixels
depending on the field type, is first handed to PaddleOCR, which is the
main OCR engine. PaddleOCR carries out a three-stage pipeline:

\begin{enumerate}
\def\labelenumi{\arabic{enumi}.}
\item
  text detection by DB (Differentiable Binarization) algorithm that
  finds the text location in the cropped image,
\item
  text recognition by CRNN (Convolutional Recurrent Neural Network) that
  changes the image of a text to a string of characters,
\item
  confidence scoring on both characters and the whole line.
\end{enumerate}

The PaddleOCR output is a list of tuples: {[}(\textit{bbox\_coords},
(\textit{text}, \textit{confidence})){]} where \textit{bbox\_coords} are the
coordinates of the text in the region, \textit{text} is the recognized
string, and \textit{confidence} is a float between 0 and 1.

In a case when PaddleOCR confidence is lower than 0.6 threshold, which
means there are almost certainly problems with recognition in the
regions that are degraded, skewed, or low-contrast, the system turns on
Tesseract OCR to a fallback mechanism automatically. Tesseract is run on
the same area

and provides a text string and a confidence score. It is equipped with
the English language model and page segmentation mode 6 (uniform block
of text). The system then pits the confidence scores from both sources
against each other and chooses the one with the higher confidence, hence
a confidence-based arbitration strategy is implemented.

There is only one \textit{product\_table} exception: this region is not
undergoing OCR now. Rather, its bounding box coordinates are recorded as
a JSON-serialized string given in the format \{"x1": float, "y1": float,
"x2": float, "y2": float\} in the \textit{product\_table\_bbox field},
with a reason flag saying "Table extraction handled by post-processing."
This clever move postpones breaking down the complicated table to the
specialized Stage 5 pipeline.

The OCR output of the rest seven fields is saved in a structured manner.
Each field is given a raw text value (e.g., \textit{gstin\_raw}) and a
confidence score (e.g., \textit{gstin\_confidence}). The machine also
stores metadata like \textit{ocr\_engine\_used} that may have values such
as ``paddleocr'',
``tesseract'', or
``paddleocr\_with\_tesseract\_fallback''
depending on the engines that were actually used. By using this
dual-engine method, the overall character recognition rates are 95.3\%
for printed text and 89.7\% for slightly degraded documents, thus the
system is leveraging single-engine implementations significantly. The
OCR stage is generally done within 1-2 seconds per region, and the whole
stage takes about 8-14 seconds for a typical 7-region document, which is
the most significant time slice of roughly 50-60\% of the total processing time.

\subsection{Raw Data Persistence}

After OCR extraction, the pipeline moves over to the database persistence
stage. Here, the raw OCR outputs are saved in PostgreSQL database with
\textit{status} as ``OCR\_COMPLETE''. This
temporary storage has three main purposes: (1) it makes the system
fault-tolerant by allowing the resumption of the processing from this
checkpoint in case the following stages crash; (2) it offers full
traceability for debugging and model improvement by keeping the original
OCR text that is even before any corrections; and (3) it saves the data
from being lost due to validation failures in post-processing.

An important architectural decision at this stage deals with database
schema constraints: the fields with strict length limitations, that are
\textit{gstin} (varchar 15) and \textit{mobile\_no} (varchar 15), will keep
the raw OCR text only in separate raw columns (\textit{gstin\_raw},
\textit{mobile\_no\_raw}), while their main columns (\textit{gstin},
\textit{mobile\_no}) will be initially NULL. This avoids
StringDataRightTruncation errors that result when OCR incorrectly reads
a non-GSTIN text as a GSTIN field, which is quite a frequent problem in
documents with variable layouts. The \textit{\_raw} columns are not
limited in length and thus can accommodate any OCR result of up to 100
characters (after truncation for safety). Thus, the entire OCR result is
kept for later analysis.

For the rest of the fields: \textit{date\_of\_receipt\_raw},
\textit{invoice\_no\_raw}, \textit{total\_amount\_raw},
\textit{store\_name\_raw}, \textit{store\_address\_raw}, the system keeps
the raw OCR text as well as the initial values, that are going to be
validated and corrected in Stage 5. Each field triplet has: (1) the
\textit{main} field where the final validated value is stored, (2) the
\textit{\_raw} field that keeps the original OCR output, (3) the
\textit{\_confidence} field that stores OCR confidence scores and (4) the
\_\textit{bbox} field that keeps the spatial coordinates in JSON format.
The \textit{full\_ocr\_text} field is a combination of all the text that
has been extracted and is there for easy reference and full-text
searching. Moreover, ensemble detection metadata like
\textit{detection\_confidence} (average across all regions) and
\textit{ensemble\_model\_used}, are also saved at this stage.

The database commit is done by using SQLModel' s
\textit{session.add()} followed by \textit{session.commit()}, and is wrapped
by a complete error handling which catches connection timeouts,
constraint violations, and transaction conflicts, thus automatically
leading to \textit{session.rollback()} and retry logic being activated on
transient failures. In the case of long-running database commits that
take more than 30 seconds due to serverless Neon PostgreSQL network
latency, the system will perform an exponential backoff retry with a
maximum of 3 retries. This persistence point normally takes from 0.5 to
1.0 seconds and about 5\% of the total processing time and is a stable
checkpoint from which the pipeline can proceed validation and correction
safely.

\subsection{Post-Processing, Validation, and LLM Correction}

This segment
of the pipeline features a hybrid validation as well as a correction
system which combines domain-specific rule-based logic with LLM-powered
semantic refinement. Any part of the document that has been processed
and the fields extracted goes through a three-step process: initial
validation, conditional LLM correction, and final re-validation.
Besides, \textit{gstin}, \textit{mobile\_no}, \textit{invoice\_no},
\textit{date}, and \textit{total\_amount} have validators especially
designed for them. Here, format-specific constraints are imposed by
these validators, for example, regex matching for GSTIN structure and
Indian state codes, extraction of 10-digit mobile numbers from noisy OCR
outputs, and date normalization to ISO 8601. In addition to this,
monetary values get effected in such a manner that their decimal
precision is kept, at the same time, currency symbols are removed, which
helps in fixing most of the common parsing mistakes that result from OCR
artifacts.

In case of unsuccessful validation or obtaining of low confidence for a
field, this field is forwarded to a fine-tuned, 4-bit quantized version
of Google's MedGemma-4B IT optimized for medical-domain
reasoning and deployed using Hugging Face Transformers with BitsAndBytes
quantization to enable local inference on limited VRAM. MedGemma gets
the prompt with customized instructions for each field type. For
example, store names get the benefit of such a cleaning by abbreviation
expansion (e.g., ``PVT'' to ``Pvt. Ltd.''), title casing is applied, and
spelling mistakes are corrected, whereas for store addresses proper
punctuation and abbreviation expansion are used for formatting.
Subsequently, LLM outputs undergo re-validation through the same
rule-based methods. The system at hand is able to assign correction
method tags ('validator',
'llm', or
'validator\_failed') and can also
calculate the final confidence level by using a weighted strategy
depending on correction origin and validation confidence.

For intricate document content especially the \textit{product\_table}
field, the system turns to an intelligent LLM-based parser rather than
using a grid detection method.
The product table extraction module is the part that deals with the problem
of extracting tabular data from borderless or irregular medical bills.
As a first step, the ensemble detection model localizes the tentative
table regions, however, the bounding boxes it produces are rough and
sometimes even contain parts of the page that are not related to the
table. To do this more precisely, the system utilizes Microsoft's Table
Transformer Detection model which employs a DETR-based architecture
trained on PubTables-1M in order to pinpoint tables very precisely even
if there are no lines visible - 94.2\% average precision is attained.
Though the Table Transformer also offers a structure recognition module,
it fails on more than 70\% of medical bills in the dataset because of
inconsistent or missing cell boundaries. Consequently, a hybrid approach
is used: table detection is refined by Table Transformer, then OCR and
LLM-driven semantic parsing are performed.

After locating the table region, PaddleOCR gets the text blocks together
with their spatial coordinates. These blocks are grouped into rows based
on the closeness of their y-coordinates, and then the order of the rows
is determined from left to right. The text-formatted data is sent to
MedGemma together with the prompt telling it to extract structured
product entries as JSON objects. MedGemma is capable of understanding
spatial and semantic patterns which help in correct parsing. The output
records are checked and loaded into the database; thus, the accuracy of
\textit{product} and \textit{amount} fields is over 90\% which is a great
improvement of grid-based methods on unstructured
tables.

The model output is handled by a fallback-aware JSON parsing logic which
is designed to handle even the malformed outputs. Valid product objects
are therefore created and saved as separate records in the database each
with positional row indices. Also, all the financial figures are
transformed into floats after the non-numeric characters have been
removed thus making the persistence clean. The table extraction
component has been able to achieve more than 90\% accuracy on the main
fields in different formats which is a great deal of performance
improvement over standard Hough-transform-based methods on unstructured
tables.

All the activities in this last stage of the pipeline: correction of the
fields, LLM refinement, and table parsing are done within 3-5 seconds
per document and constitute about one-quarter of the entire pipeline's
runtime. So as to facilitate audit and optimization, metrics such as
corrected fields, LLM invocation counts, and extracted product rows are
recorded. To ensure data consistency, the mechanism for the handling of
transactions is very robust, and in cases of transient database errors,
it is equipped with automatic rollback and retry logic, thus, the stage
of correction being intelligent, sturdy, and ready for large-scale
deployment is brought to a
close.

\subsection{Finalization and Status Update}

The last stage consolidates the status, gathers all errors, and prepares the
response. If everything has gone well in the previous stages, then the
medical document status is locally changed from
``OCR\_COMPLETE'' to
``COMPLETED'' by means of a database
transaction, with \textit{processing\_timestamp} reflecting the current
UTC time. The system builds up a very detailed JSON payload comprising:
(1) the full \textit{medicalBill} model containing all the extracted and
verified data fields, (2) a \textit{detection\_summary} object holding the
ensemble predictions for each of the eight classes along with bbox
coordinates, class labels, and confidence scores, (3) an
\textit{ocr\_results} object associating the field names with raw OCR
text, engine used, and confidence scores, (4) a \textit{post\_processing}
object having \textit{fields\_corrected},
\textit{product\_items\_extracted}, and \textit{llm\_corrections} counts and
(5) a \textit{product\_items} list encompassing all the extracted table
rows, each of them being \{\textit{product}, \textit{quantity}, \textit{pack},
\textit{mrp}, \textit{expiry}, \textit{total\_amount}\} formatted. The server
response does not contain the raw binary data nor the internal IDs.

When an exception is raised in any stage, the pipeline reacts by
graceful degradation: (1) the failed stage \textit{ProcessingJob} is
updated to \textit{status} with ``FAILED''
along with detailed \textit{error\_message} capturing the exception
traceback, (2) the corresponding \textit{MedicalBill} \textit{status} is
changed to ``FAILED'' with
\textit{validation\_notes} holding a short failure description, (3)
\textit{session.rollback()} is executed so the partial commits are
avoided, (4) the error response with HTTP status code 500 and a JSON
body containing the \textit{error\_message} is delivered back to the
client side to inform it about the exact point of failure. This
finalization stage, takes about 0.3-0.5 seconds, hence, it is
responsible for less than 3\% of the overall processing time.

\section{Results and Performance Analysis}

\subsection{End-to-End Pipeline Performance}

The pipeline was tested on the test set of 60 medical bills that represented
various layouts, scan qualities, and document conditions. The system
exhibited excellent performance with a success rate of 93.3\% at the
document level. On average the pipeline took around 17.3 seconds to
process a single document. The performances were measured on my system
equipped with an NVIDIA RTX 3050 GPU (6GB VRAM), AMD Ryzen™ 5 7235HS
processor, and 24GB RAM. Time durations measured for individual stages
are as follows: initialization (0.3s), ensemble detection (2.7s),
dual-engine OCR (9.8s), database persistence (0.9s), post-processing and
validation (3.2s), and finalization (0.4s). The OCR extraction part was
approximately 57\% of the whole operation time. It was followed by
post-processing that took 18\% and ensemble detection that took 16\%.
This means that text recognition is still the main computational
bottleneck although GPU acceleration is
used.

\subsection{Field Level Extraction Accuracy}

Each field's extraction accuracy was locally improved per
eight target classes by measuring exact match criteria for structured
fields and semantic similarity scoring for free-text fields. The highest
accuracy has been achieved on total\_amount (96.7\%), invoice\_no
(95.0\%), and date\_of\_receipt (93.3\%) that were helped by strong
regex-based validators and the consistent formatting patterns of
documents. Store\_name and store\_address reached 91.7\% and 89.2\%
accuracy correspondingly, with the LLM-based correction that has been
successful in 78\% of initially failed validations by abbreviation
expansion, title casing, and spelling correction. Store\_address
continued with 90.0\% accuracy, with details of the failure being
OCR-induced character confusion in "GP" instead of "OP" in "Murugan GP
Market Road Indore" and recognition errors of extracted tokens that led
to the words "Murugan" and "Market" being unlinked in the processed text
thus lowering overall scoring. GSTIN demonstrated 87.5\% accuracy with
major sources of failure being OCR misrecognition of similarly-looking
characters (8 vs B, 0 vs O) in the 15-character alphanumeric code even
though Luhn algorithm validation has been implemented. Mobile number
extraction turned out to be the most difficult part leading to 82.1\%
accuracy as a result of variable formatting (spaces, hyphens,
parentheses), the presence of multiple phone numbers in close proximity,
and the occasional confusions with order IDs or invoice
numbers.

\begin{table}[htbp]
\caption{Per-Class Detection Performance at IoU 0.50 and 0.75}
\label{tab:perclass}
\centering
\scriptsize
\begin{tabular}{lccccc}
\toprule
Class & IoU & Prec. & Rec. & F1 & Avg IoU \\
\midrule
\multirow{2}{*}{date\_of\_reciept} & 0.50 & 0.950 & 0.983 & 0.966 & 0.834 \\
& 0.75 & 0.817 & 0.845 & 0.831 & 0.850 \\
\multirow{2}{*}{gstin} & 0.50 & 0.870 & 0.825 & 0.847 & 0.827 \\
& 0.75 & 0.759 & 0.719 & 0.739 & 0.850 \\
\multirow{2}{*}{invoice\_no} & 0.50 & 0.918 & 0.966 & 0.941 & 0.814 \\
& 0.75 & 0.738 & 0.776 & 0.756 & 0.835 \\
\multirow{2}{*}{mobile\_no} & 0.50 & 0.618 & 0.737 & 0.672 & 0.793 \\
& 0.75 & 0.456 & 0.544 & 0.496 & 0.841 \\
\multirow{2}{*}{product\_table} & 0.50 & 0.944 & 0.985 & 0.964 & 0.900 \\
& 0.75 & 0.916 & 0.956 & 0.935 & 0.907 \\
\multirow{2}{*}{store\_address} & 0.50 & 0.891 & 0.983 & 0.934 & 0.847 \\
& 0.75 & 0.766 & 0.845 & 0.803 & 0.873 \\
\multirow{2}{*}{store\_name} & 0.50 & 0.971 & 0.971 & 0.971 & 0.855 \\
& 0.75 & 0.868 & 0.868 & 0.868 & 0.871 \\
\multirow{2}{*}{total\_amount} & 0.50 & 0.934 & 0.983 & 0.958 & 0.831 \\
& 0.75 & 0.820 & 0.862 & 0.840 & 0.848 \\
\bottomrule
\end{tabular}
\end{table}

\subsection{Product Table Extraction Results}

A hybrid table extraction strategy that integrates Table Transformer
detection, PaddleOCR text localization, and MedGemma semantic parsing
has led to a 91.2\% accuracy level in structured product rows. This was
based on full row extraction where all six fields (product name,
quantity, pack size, MRP, expiry date, and total amount) were correctly
identified. Of the 486 product rows in the test set, 443 were perfectly
extracted, 28 had one or two fields either missing or incorrect (partial
extraction), and 15 were complete failures due to highly irregular table
layouts that did not have any visual structure.

Borderless tables (68\% of test cases), multi-line product names (34\%
of cases), and merged cells (22\% of cases) were successfully dealt with
by the system, thus it has demonstrated significant superiority to
traditional grid-detection methods which have only 64\% accuracy on the
same dataset. The per-field extraction accuracy of product tables was:
\textit{product\_name} (94.2\%), \textit{quantity} (89.7\%),
\textit{pack\_size} (88.1\%), \textit{mrp} (92.5\%), \textit{expiry\_date}
(86.3\%), and \textit{amount} (93.8\%), with date fields having the
highest error rate due to different formatting of dates by
manufacturers.

\subsection{LLM Correction Impact Analysis}

The post-processing step with MedGemma integration was a major factor in the
improvement of the overall accuracy, as it made corrections on 42.3\% of
the total extracted fields (256 out of 605 field instances across 60
documents). The rate of LLM intervention was 83.6\%, as the system was
able to make the correct changes in 214 fields that initially failed
rule-based validation, while 42 fields were left for manual review even
after LLM processing. The most frequent types of correction operations
were: date format normalization (changing DD/MM/YYYY, MM-DD-YY, and
different non-standard formats to ISO 8601), which accounted for 31\% of
the corrections; numeric cleaning (removing currency symbols, commas,
and characters that are not necessary from monetary values), which
accounted for 24\%; store name standardization (abbreviation expansion,
title casing, and spelling correction), which accounted for 19\%;
address formatting (punctuation correction, landmark expansion and
standardization), which accounted for 15\%; and GSTIN/mobile number
cleaning (removing spaces, hyphens, and characters that are not valid),
which accounted for 11\%. The average LLM inference time was 0.8
seconds/field on quantized MedGemma 4B, and 4-bit quantization was able
to reduce VRAM requirements from 16GB to 3.8GB while correction accuracy
was kept within 2\% of the full-precision
model.

\subsection{Comparative Analysis with Baseline Models}

The ensemble-based pipeline put up against three baseline methods on the
identical test set: (1) solo run of YOLOv11n with PaddleOCR, (2)
standalone Faster R-CNN with PaddleOCR, and (3) template-based rule
extraction using Tesseract OCR; yielded a 91.4\% accuracy at the field
level, which is higher than 87.9\% achieved by YOLOv11n alone, 84.3\% by
Faster R-CNN alone, and 68.7\% by template-based extraction, thus
reflecting the accuracy that has been lifted by 3.5\%, 7.1\%, and 22.7\%
respectively.

On comparing processing times, the trade-offs were quite intriguing: the
ensemble system took an average of 17.3s per document, while YOLOv11n
alone took 12.8s (25\% faster but less accurate), Faster R-CNN alone
took 15.4s (11\% faster) and template-based method took only 8.2s (52\%
faster but very inaccurate and completely failing on non-standard
layouts). The F1-score study at IoU threshold 0.5 showed that WBF
ensemble (90.80\%) was better than individual models in all the fields
except in \textit{mobile\_no}, where the precision-focused approach of
Faster R-CNN just showed marginal advantages.

It is important to note that the dual-OCR method with
PaddleOCR-Tesseract as a fallback drastically improved the
character-level accuracy by 7.2\%, especially in the cases of the
documents with low-contrast or degraded regions, where confidence-based
arbitration selected the output of the engine that was
superior.

\section{Limitations and Future Work}

\subsection{Current System Limitations}

Even though the pipeline manages to deliver powerful results on the test
dataset, several limitations affect the strength and the range of the
system. The first limitation is that the dependency on bounding box
annotations during the training phase means that the models will not be
able to find fields that differ a lot from the eight classes that have
been predefined; this means that the adaptability to bills containing
additional information types such as tax breakdowns, discount codes, or
patient details will be very limited. \cite{ref2}

Secondly, the pipeline performance drops very significantly on documents
that are heavily degraded - images with resolution less than 600×800
pixels, blur more than motion kernel size 15 or skew angles beyond ±25°
result in failure rates more than 40\% because the preprocessing module
does not have advanced document rectification algorithms \cite{ref2},
\cite{ref5}.

Thirdly, the mobile number detection at 82.1\% reflects that the main
OCR challenge is in differentiating between the very similar digit
patterns; this problem gets worse when several phone numbers are written
close to each other without any labels and characters. Moreover, object
detection models like YOLO, on which this pipeline relies, have been
known to struggle with predicting small objects that appear in groups
\cite{ref5}, \cite{ref15}.

Fourthly, the single-threaded processing architecture which is designed
to process one document at a time only, is limiting the throughput to
about 208 documents per hour, therefore, the GPU parallelization
opportunities are not fully utilized \cite{ref3}, \cite{ref4}, \cite{ref17}.

Fifthly, the system is currently supporting only English and numerical
text, thus it is failing in cases of Hindi or any other regional
language content which makes up for about 15-20\% of medical bills in
rural areas of India, as the English model of PaddleOCR and the
English-focused training of MedGemma restrict that the system is not
able to handle multiple languages \cite{ref4}, \cite{ref5}, \cite{ref10}, \cite{ref19}.

\subsection{Proposed Enhancements}

Several upgrades across the board for preprocessing, modeling, data quality, and
validation layers are planned to extend the capabilities of the proposed
system. One of the ways the system could be enhanced is by using
different methods for image preprocessing. For example, for document
deskewing \cite{ref2}, text detection based on CRAFT \cite{ref7} could be used;
for super-resolution, ESRGAN could be utilized; and for noise
suppression, adaptive bilateral filtering could be applied-these are all
techniques that are expected to lead to an OCR accuracy improvement of
8-12\% on degraded scans. By gradually training the model with new data,
the detection model can be extended to support 15-20 classes and thus it
would be possible to get the values of auxiliary fields such as batch
numbers and manufacturer details \cite{ref2}. Hindi, Tamil, and other
scripts in India could be supported if the system were to be extended
multilingual-wise using IndicOCR or Bhashini \cite{ref5}, thereby the system
would be available for the whole country.

Consideration of asynchronous task and containerized horizontal scaling
by means of Docker and Kubernetes may result in the
system's throughput to be elevated to the level of
thousands of documents per hour. Improvement of data quality is still a
challenging problem as the existing dataset consisting of 862 images of
documents limits the generalizability of the model. 2,000-3,000
annotated bills can be obtained through crowd-sourcing \cite{ref4}, \cite{ref5},
synthetic augmentation, and semi-supervised learning to scale the
dataset which will not only increase the robustness of the model but
also the annotation quality via active learning, consensus labeling, and
automated quality checks \cite{ref3}.

Moreover, the solution to the problem of class imbalance, especially in
the case of underrepresented entities such as \textit{mobile\_no} could be
the use of class-weighted loss functions and targeted collection for
performance improvement of minority classes by identifying the most
suitable data sampling strategies. The error hunting process may be
supported by machine learning-based anomaly detection as well as
cross-field consistency checks (e.g., table total matching overall
amount) while done from the side of validation. Machine learning
techniques can be used for anomaly detection and cross-field consistency
checks (e.g., table total matching overall amount) to identify error
sources better and more quickly \cite{ref9}. Selective automations on the
basis of confidence and active learning feedback loops are just two of
the ways human-in-the-loop correction can be supported, thus allowing
for continuous improvement facilitated by the balancing of accuracy and
scalability for deployment in the real
world.

\subsection{Research Directions}

Several promising research directions are identifiable from this work. One
possibility is that attention-based ensemble methods that dynamically
weight model predictions based on input characteristics (e.g., document
quality metrics, field complexity scores) instead of static
class-specific weights could lead to better adaptive performance.
Another angle to consider is the effect of multimodal fusion
architectures that jointly process the visual features (raw image
patches) and the textual features (OCR outputs) through
transformer-based encoders, which might surpass the sequential
visual-then-textual processing in terms of semantic understanding
\cite{ref5}, \cite{ref10}. A further idea is the creation of domain-specific
vision-language models by continued pre-training of foundation models
like Florence-2 or Qwen-VL on large-scale medical document corpora that
would result in specialized extractors becoming far more efficient than
general-purpose OCR + LLM pipelines \cite{ref3}, \cite{ref11}, \cite{ref19}. Apart
from that, the study of zero-shot and few-shot learning methods that
allow quick adaptation to new bill formats with a few labeled examples,
thus, the annotation demand for healthcare providers diversification,
would be considerably lowered, is also a viable direction. The last
point is the usage of differential privacy and federated learning
methods to facilitate collaborative model enhancement across healthcare
institutions without the exchange of sensitive patient data, which is a
significant step towards privacy-preserving system
improvement.

\section{Conclusion}

This paper discusses the design, implementation, and evaluation of an
end-to-end document parsing pipeline for automatically extracting and
validating structured data from Indian medical bills. The pipeline
integrates multiple cutting-edge techniques: a weighted box fusion
ensemble combining YOLOv11n and Faster R-CNN for precise field
detection, a dual-engine OCR approach with PaddleOCR and Tesseract for
robust text extraction, Microsoft's Table Transformer
for identifying complex tabular regions, and MedGemma 4B, a
domain-specialized model, for intelligent semantic correction and
validation. On a self-annotated test set of 60 varied medical bills, the
pipeline achieved a 93.3\% document-level success rate and a 91.4\%
field-level extraction accuracy, with an average processing time of 17.3
seconds per document on consumer-grade GPU hardware.

The ensemble detection method was very efficient and managed to achieve
a macro F1-score of 90.66\% at IoU 0.5 by utilizing class-specific
fusion weights, which allowed the optimization of the complementary
strengths of real-time YOLO and precision-focused Faster R-CNN
architectures. The dual-OCR strategy with confidence-based arbitration
result character-level accuracy improvements by 7.2\% over single-engine
baselines, thus proving the necessity of redundancy in dealing with
variable document quality. The LLM-based post-processing pipeline
managed to fix 83.6\% of the initial errors substantially lowering the
requirement for manual review, at the same time, it preserved the
semantic correctness due to the domain-aware reasoning. The hybrid table
extraction technique: employing detection, spatial OCR clustering and
LLM-driven semantic parsing reached 91.2\% accuracy in structured
product rows even though grid structures were missing in 68\% of the
cases thus it was 27\% better than the traditional grid-detection
method.

The modular design, which clearly separates detection, extraction,
validation, and correction stages, makes it possible to gradually bring
in improvements and allows the system to be put in place on affordable
local hardware via 4-bit model quantization and eager loading
strategies. The pipeline is capable of processing roughly 208 documents
per hour with a single RTX 3050 GPU and can be easily scaled
horizontally for institutional-scale workloads. Although the existing
limitations, such as the system's sensitivity to heavily
degraded images, the absence of multilingual support, and restrictions
on the extensibility of field types, are still challenges to be
addressed later on, the pipeline is proof that carefully coordinated
combinations of specialized models can produce strong document
comprehension without the need for large-scale end-to-end training on
domain-specific data.

This research is a step forward in the AI industry-specific deployment
literature for healthcare automation. It proves that hybrid methods that
fuse classical computer vision, state-of-the-art deep learning, and
large language models can bring about production-ready systems that
substantially enhance the throughput of manual work yet still maintain
high accuracy standards that are compatible with real-world medical
billing workflows. The solutions and architectural strategies put
forward here have a wider application than medical bills alone,
extending to any structured document extraction tasks and thus offering
a framework for the construction of dependable, serviceable, and
high-performing document intelligence systems. Planned future research
will concentrate on widening the range of languages supported, enhancing
the degraded document performance through better preprocessing, and
investigating the possibility of using end-to-end vision-language models
that not only simplify the pipeline but also keep or exceed the current accuracy levels.

\section*{Acknowledgment}

  I would like to express my sincere gratitude to my professor Dr. Kanika
Garg, for mentoring me throughout the duration of this research. I also
thank the SRMIST Delhi-NCR Campus for providing resources and
institutional support that made this research possible. Special thanks
to the open-source community for developing and maintaining the
foundational libraries and models utilized in this work. I am grateful
to the Roboflow Universe platform for hosting the initial dataset and to
all contributors who make their work publicly available, advancing
research and development in document understanding and medical AI.

\balance
\begin{thebibliography}{20}

\bibitem{ref1}
R. L. Chaudhary, A. R. Peddireddy, and U. Mathur, ``Robotic Process Automation in Billing Management: Enhancing Healthcare Efficiency,'' 2023, \textit{Unpublished}. doi: 10.13140/RG.2.2.27402.68809.

\bibitem{ref2}
Š. Šimsa \textit{et al.}, ``DocILE Benchmark for Document Information Localization and Extraction,'' May 2023, \textit{arXiv}: arXiv:2302.05658. doi: 10.48550/arXiv.2302.05658.

\bibitem{ref3}
M. Dhouib, G. Bettaieb, and A. Shabou, ``DocParser: End-to-end OCR-free Information Extraction from Visually Rich Documents,'' May 2023, \textit{arXiv}: arXiv:2304.12484. doi: 10.48550/arXiv.2304.12484.

\bibitem{ref4}
G. Kim \textit{et al.}, ``OCR-Free Document Understanding Transformer,'' in \textit{Computer Vision -- ECCV 2022}, vol. 13688, Cham: Springer Nature Switzerland, 2022, pp. 498--517. doi: 10.1007/978-3-031-19815-1\_29.

\bibitem{ref5}
S. Bhattacharyya, P. Deb, and C. V. Jawahar, ``Adapting Vision-Language Models for Hindi OCR''.

\bibitem{ref6}
A. Datta \textit{et al.}, ``Preserving medical information from doctor's prescription ensuring relation among the terminology,'' \textit{Comput. Biol. Med.}, vol. 188, p. 109812, Apr. 2025. doi: 10.1016/j.compbiomed.2025.109812.

\bibitem{ref7}
Y. Baek, B. Lee, D. Han, S. Yun, and H. Lee, ``Character Region Awareness for Text Detection,'' in \textit{2019 IEEE/CVF CVPR}, Long Beach, CA, USA, June 2019, pp. 9357--9366. doi: 10.1109/CVPR.2019.00959.

\bibitem{ref8}
D. Deng, H. Liu, X. Li, and D. Cai, ``PixelLink: Detecting Scene Text via Instance Segmentation,'' Jan. 2018, \textit{arXiv}: arXiv:1801.01315. doi: 10.48550/arXiv.1801.01315.

\bibitem{ref9}
Z. Shen, R. Zhang, M. Dell, B. C. G. Lee, J. Carlson, and W. Li, ``LayoutParser: A Unified Toolkit for Deep Learning Based Document Image Analysis,'' June 2021, \textit{arXiv}: arXiv:2103.15348. doi: 10.48550/arXiv.2103.15348.

\bibitem{ref10}
Y. Huang, T. Lv, L. Cui, Y. Lu, and F. Wei, ``LayoutLMv3: Pre-training for Document AI with Unified Text and Image Masking,'' July 2022, \textit{arXiv}: arXiv:2204.08387. doi: 10.48550/arXiv.2204.08387.

\bibitem{ref11}
B. Davis, B. Morse, B. Price, C. Tensmeyer, C. Wigington, and V. Morariu, ``End-to-end Document Recognition and Understanding with Dessurt,'' June 2022, \textit{arXiv}: arXiv:2203.16618. doi: 10.48550/arXiv.2203.16618.

\bibitem{ref12}
X.-H. Li, F. Yin, and C.-L. Liu, ``DocSAM: Unified Document Image Segmentation via Query Decomposition and Heterogeneous Mixed Learning''.

\bibitem{ref13}
S. Ren, K. He, R. Girshick, and J. Sun, ``Faster R-CNN: Towards Real-Time Object Detection with Region Proposal Networks,'' Jan. 2016, \textit{arXiv}: arXiv:1506.01497. doi: 10.48550/arXiv.1506.01497.

\bibitem{ref14}
K. He, G. Gkioxari, P. Dollár, and R. Girshick, ``Mask R-CNN,'' Jan. 2018, \textit{arXiv}: arXiv:1703.06870. doi: 10.48550/arXiv.1703.06870.

\bibitem{ref15}
J. Redmon, S. Divvala, R. Girshick, and A. Farhadi, ``You Only Look Once: Unified, Real-Time Object Detection,'' in \textit{2016 IEEE CVPR}, Las Vegas, NV, USA, June 2016, pp. 779--788. doi: 10.1109/CVPR.2016.91.

\bibitem{ref16}
N. Livathinos \textit{et al.}, ``Docling: An Efficient Open-Source Toolkit for AI-driven Document Conversion,'' Jan. 2025, \textit{arXiv}: arXiv:2501.17887. doi: 10.48550/arXiv.2501.17887.

\bibitem{ref17}
K. Lee \textit{et al.}, ``Pix2Struct: Screenshot Parsing as Pretraining for Visual Language Understanding,'' June 2023, \textit{arXiv}: arXiv:2210.03347. doi: 10.48550/arXiv.2210.03347.

\bibitem{ref18}
D. Lewis, G. Agam, S. Argamon, O. Frieder, D. Grossman, and J. Heard, ``Building a test collection for complex document information processing,'' in \textit{Proc. 29th ACM SIGIR}, Seattle, WA, Aug. 2006, pp. 665--666. doi: 10.1145/1148170.1148307.

\bibitem{ref19}
A. Sellergren \textit{et al.}, ``MedGemma Technical Report,'' July 2025, \textit{arXiv}: arXiv:2507.05201. doi: 10.48550/arXiv.2507.05201.

\bibitem{ref20}
W. Kwon \textit{et al.}, ``Efficient Memory Management for Large Language Model Serving with PagedAttention,'' Sept. 2023, \textit{arXiv}: arXiv:2309.06180. doi: 10.48550/arXiv.2309.06180.

\end{thebibliography}

\end{document}



